\textbf{Background:} Information-gain utilities make sensing valuable directly, but introduce a weight that is usually tuned per task. An engineer may instead need to specify expected sensing usage.
\textbf{Objectives:} We calibrate the information weight of a $\rho$-POMDP policy family to a sensing target and identify the role of Expected Free Energy (EFE) within that family.
\textbf{Methods:} We define an operational shadow price as a crossing threshold of the usage curve, distinguish it from a finite-grid bracket, and construct an endpoint mixture for targets inside usage gaps. We prove scale equivariance of the receding-horizon planner and stationary convergence of a projected controller under stated conditions. Minimizing the specified recursive, reward-based EFE objective is equivalent to maximizing reward plus information gain at $w=1$ for observe-then-commit problems and hidden-state-preserving actions in a factored extension. The reward--information identification is exact under log scoring and otherwise depends on the declared reward convention.
\textbf{Results:} Usage curves collapse after rescaling, and the controller re-adapts empirically after a tenfold reward shift. Targets selected by a predeclared rule from calibration curves are met on held-out seeds to within $0.13$ observations. Calibration can nevertheless sacrifice reward relative to spending less or choosing another mixture. Against a sampled constrained reference, nominal pointwise paired intervals show reward shortfalls on eight of eleven distinct budgets. These intervals omit reference-selection uncertainty, and two Tiger comparisons are sensitive to rare errors absent from the held-out episodes. The untuned $w=1$ lies in the reward-maximizing bracket on three of five swept environments and is beaten on the other two.
\textbf{Conclusions:} Usage calibration is appropriate when expected sensing usage is the design requirement. It does not maximize reward under a cap. EFE supplies an information-unit anchor under specified assumptions, while usage, reward, and success must be evaluated separately.
